\documentclass[11pt,a4paper]{article}
\usepackage{fontspec}
\usepackage{polyglossia}
\setdefaultlanguage{russian}
\setmainfont{PT Serif}
\setsansfont{PT Sans}
\setmonofont{PT Mono}
\newfontfamily\cyrillicfont{PT Serif}
\newfontfamily\cyrillicfonttt{PT Mono}
\usepackage[margin=25mm]{geometry}
\usepackage{tabularx,longtable,fancyvrb}
\usepackage[hidelinks]{hyperref}
\setlength{\parindent}{1.25em}
\setlength{\emergencystretch}{3em}
\setlength{\parskip}{0.3em}
\renewcommand{\arraystretch}{1.25}
\hypersetup{pdftitle={ИЫ.00110-01 51 01},pdfsubject={ВЫЧИСЛИТЕЛЬНЫЙ КОМПЛЕКС "ЭЛЬБРУС 1-КБ" ОБЩЕЕ СИСТЕМНОЕ ПРОГРАММНОЕ ОБЕСПЕЧЕНИЕ ПОРЯДОК И МЕТОДИКА ИСПЫТАНИЙ}}
\begin{document}
\begin{titlepage}\small
\noindent\begin{minipage}[t]{0.45\linewidth}\end{minipage}\hfill\begin{minipage}[t]{0.45\linewidth}\raggedleft \end{minipage}
\vspace*{\fill}\begin{center}
{\large Вычислительный комплекс «Эльбрус 1-КБ»}\\[0.6em]
{\large Общее системное программное обеспечение}\\[0.6em]
{\large Порядок и методика испытаний}\\[0.6em]
\vspace{1em}{\large\bfseries ИЫ.00110-01 51 01}\\[0.4em]
Листов 56\end{center}\vspace*{\fill}
\begin{flushright}\begin{tabular}{r}\end{tabular}\end{flushright}
\begin{center}1988\end{center}\end{titlepage}
\tableofcontents\clearpage
\section*{1.~Объект испытаний}\addcontentsline{toc}{section}{1.~Объект испытаний}
Объект испытаний --- Общее системное программное обеспечение (ОСПО) вычислительного комплекса «Эльбрус 1-КБ». Обозначение ОСПО --- ИЫ.00110-01. ОСПО предназначено для организации вычислительного процесса на ВК «Эльбрус 1-КБ» (ВК Э1-КБ).

\section*{2.~Цель испытаний}\addcontentsline{toc}{section}{2.~Цель испытаний}
Испытания проводятся с целью:

\par\textbf{2.1.} Проверки правильности функционирования общего системного программного обеспечения на ВК «Эльбрус 1-КБ» и оценки комплектности программной документации;\par
\par\textbf{2.2.} Определения пригодности ОСПО для передачи в эксплуатацию на ВК «Эльбрус 1-КБ».\par
\section*{3.~Требования к программе}\addcontentsline{toc}{section}{3.~Требования к программе}
\subsection*{3.1.~Требования по составу ОСПО}\addcontentsline{toc}{subsection}{3.1.~Требования по составу ОСПО}
В состав общего системного программного обеспечения включаются:

\par\hangindent=1.5em - Операционная система (ОС) Диспак;\par
\par\hangindent=1.5em - Системы программирования режима совместимости (системы программирования ЭВМ БЭСМ-6 и МВК «Эльбрус 1-К2»);\par
\par\hangindent=1.5em - Система программирования ВК «Эльбрус 1-КБ», позволяющая работать как в режиме короткого адреса (режим 1), так и в режиме расширенного формата (режим 3) (см. «ВК Эльбрус 1-КБ. Техническое описание» ИЫ3.055.045 ТО5);\par
\par\hangindent=1.5em - Сервисные программы;\par
\par\hangindent=1.5em - Комплект тестовых задач для проверки работоспособности вычислительной системы;\par
\par\hangindent=1.5em - Тесты для проверки работоспособности оборудования;\par
\par\hangindent=1.5em - Комплект программной документации.\par
\begin{center}(п.~6.3)\end{center}
\subsection*{3.2.~Требования к техническим характеристикам}\addcontentsline{toc}{subsection}{3.2.~Требования к техническим характеристикам}
ОСПО должно обеспечивать:

\par\textbf{3.2.1.} Непрерывную мультипрограммную работу вычислительной системы в режимах пакетной обработки, дистанционно-пакетной обработки, работу пользователей с терминалов в режиме разделения времени, совмещение указанных режимов.\par
\begin{center}(п.~6.17)\end{center}
\par\textbf{3.2.2.} Формирование и накопление задач во входной очереди до 255.\par
\begin{center}(п.~6.16)\end{center}
\par\textbf{3.2.3.} Возможность управления планированием мультипрограммной смеси задач в соответствии с потребностями вычислительного центра.\par
\begin{center}(п.~6.14)\end{center}
\par\textbf{3.2.4.} Ведение и контроль бюджетно-статистического учета использования ресурсов.\par
\begin{center}(п.~6.8, 6.15)\end{center}
\par\textbf{3.2.5.} Вывод информации на устройства вывода.\par
\begin{center}(п.~6.6)\end{center}
\par\textbf{3.2.6.} Обработку аварийных ситуаций, возникающих при работе внешних устройств.\par
\begin{center}(п.~6.7)\end{center}
\par\textbf{3.2.7.} Возможность вывода информации на конкретное устройство вывода, на свободное устройство, на выделенное устройство.\par
\begin{center}(п.~6.6)\end{center}
\par\textbf{3.2.8.} Возможность использования архивной системы для распределения пространства на магнитных носителях.\par
\begin{center}(п.~6.4)\end{center}
\par\textbf{3.2.9.} Средства бюджетирования архивного пространства и защиту информации пользователя от несанкционированного доступа.\par
\begin{center}(п.~6.4, 6.5)\end{center}
\par\textbf{3.2.10.} Возможность пакетной и интерактивной отладки программ, написанных на языках высокого уровня в символьном и кодовом режимах.\par
\begin{center}(п.~6.13)\end{center}
\par\textbf{3.2.11.} Разметку магнитных дисков с отбраковкой дефектных участков.\par
\begin{center}(п.~6.2)\end{center}
\par\textbf{3.2.12.} В режиме 3 должно обеспечиваться решение задач пользователя, написанных на ассемблере «Эльбрус 1-КБ».\par
\begin{center}(п.~6.9, 6.10, 6.11)\end{center}
\par\textbf{3.2.13.} В режиме 3 должно обеспечиваться решение задач пользователя, написанных на фортране (транслятор FTN-В).\par
\begin{center}(п.~6.12)\end{center}
\section*{4.~Требования к программной документации}\addcontentsline{toc}{section}{4.~Требования к программной документации}
\subsection*{4.1.~Состав предъявляемой документации}\addcontentsline{toc}{subsection}{4.1.~Состав предъявляемой документации}
При проведении испытаний предъявляется документация на ОСПО согласно ведомости эд ИЫ.00110-01 20 01 и формуляру ИЫ.00110-01 30 01.

\begin{center}(п.~6.1)\end{center}
\subsection*{4.2.~Требования к предъявляемой документации}\addcontentsline{toc}{subsection}{4.2.~Требования к предъявляемой документации}
Документация на ОСПО должна быть комплектна, написана и оформлена в соответствии с требованиями гостов ЕСПД.

\begin{center}(п.~6.1)\end{center}
Примечание: в скобках приведены пункты методики испытаний, проверяющие данные требования.

\section*{5.~Средства и порядок испытаний}\addcontentsline{toc}{section}{5.~Средства и порядок испытаний}
\subsection*{5.1.~Технические и программные средства}\addcontentsline{toc}{subsection}{5.1.~Технические и программные средства}
\par\textbf{5.1.1.} Аппаратура комплекса «Эльбрус 1-КБ» должна соответствовать стандартной конфигурации ВК «Эльбрус 1-КБ» и правильно функционировать на тестово-диагностических программах из состава СТП ИЫ.04301-01.\par
На пультовой ЭВМ должно правильно функционировать пультовое программное обеспечение (ППО) (см. «ППО. Руководство оператора» ИЫ.00112-01 34 01).

\par\textbf{5.1.2.} Документация на ОСПО должна быть распечатана и предъявлена на бумаге.\par
\par\textbf{5.1.3.} Необходимые для испытаний магнитные носители (магнитные диски, магнитные ленты) должны быть подготовлены заблаговременно. Перед вызовом ОС магнитные диски (МД) должны быть размечены (см. подраздел 6.2).\par
Для испытания операционной системы должен быть предварительно сформирован магнитный диск 2048 --- Основной диск архивной системы с системой программирования ВК «Эльбрус 1-КБ», системами программирования режима совместимости, сервисными программами.

2053 --- Диск с готовыми программами ОС --- Формируется в процессе испытаний.

Оба диска являются архивными и должны быть доступны для записи.

Для демонстрации первичного вызова ОС и формирования МД 2053 должна быть подготовлена МЛ1 с готовыми программами ОСПО (см. «ОС Диспак. Руководство системного программиста» ИЫ.00110-01 32 01).

\par\textbf{5.1.4.} Для работы программ статистики должна быть заранее создана область накопления статистики по задачам и устройствам «Стат» длиной 1744 зоны. В область должны быть переписаны программы статистики и сделаны начальные генерации программ «Сбор» и «Итог» (см. «ОС Диспак. Руководство системного программиста»).\par
\par\textbf{5.1.5.} Для решения контрольной задачи КЗ3 используются четыре области длиной по 200 зон с именами «КЗ31», «КЗ32», «КЗ33», «КЗ34».\par
Эти области создаются в процессе испытаний.

\subsection*{5.2.~Последовательность испытаний}\addcontentsline{toc}{subsection}{5.2.~Последовательность испытаний}
\par\textbf{5.2.1.} Проверка комплектности программной документации производится путем сличения предъявленной документации с содержанием ведомости эд ИЫ.00110-01 20 01 и формуляра ИЫ.00110-01 30 01.\par
\par\textbf{5.2.2.} На испытания представляются готовые программы ОСПО на магнитном диске 2048 --- Основной диск архивной системы с системой программирования ВК «Эльбрус 1-КБ», системами программирования режима совместимости, сервисными программами; и на магнитной ленте МЛ1 --- Программа переписи и готовые программы операционной системы.\par
Магнитный диск 2053 формируется в процессе испытаний.

\par\textbf{5.2.3.} Перед вызовом ОС проводится демонстрация разметки магнитного диска с помощью автономной программы разметки из состава программ «Сервистст» (см. подраздел 6.2).\par
На размеченный диск переписываются готовые программы ОС с МЛ1 (см. «ОС Диспак. Руководство системного программиста» ИЫ.00110-01 32 01).

\par\textbf{5.2.4.} Производится генерация (вызов) операционной системы Диспак (см. подраздел 6.3).\par
После вызова ОС на операторский терминал с логическим номером 1 выдается сообщение об имеющихся ресурсах системы. На операторском терминале набираются директивы задания смены операторов «Сме 1», даты «Дат XXZZ» и астрономического времени «Вре ЧЧММ» (см. «Операционная система Диспак. Руководство оператора» ИЫ.00110-01 34 01), после чего система начинает работать в стандартном режиме.

Для изучения структуры системы выдается сводка по архивным томам 2048 и 2053 (см. подраздел 6.3).

\par\textbf{5.2.5.} Для проверки возможности работы с архивной системой с помощью диалоговой программы «Арфа» создаются:\par
\par\hangindent=1.5em - Четыре области с именами «КЗ31», «КЗ32», «КЗ33», «КЗ34» для решения контрольной задачи КЗ3.\par
\par\textbf{5.2.6.} Для обращения к библиотекам с контрольными задачами используется диалоговая программа «Димип», позволяющая как формировать отдельные задачи из библиотеки, так и запускать пакет контрольных задач целиком (см. «Операционная система Диспак. Руководство программиста» ИЫ.00110-01 33 01, раздел 6).\par
Программа «Димип» вызывается по ключу «Димип».

\par\textbf{5.2.7.} Проверка защиты информации в архиве от несанкционированного доступа осуществляется с помощью программы «Серп». Делается попытка чтения из чужой области. Программа «Серп» снимается с диагностикой «Чужая информация» на АЦПУ.\par
\par\textbf{5.2.8.} Возможность распределения выдач на АЦПУ осуществляется с помощью программы «Серп». Одно из АЦПУ выделяется шифру 1317.. И формируется задача с указанным шифром. Выдача данной задачи попадает на распределенное АЦПУ, выдачи других задач на это АЦПУ не попадают.\par
\par\textbf{5.2.9.} Бюджетно-статистический контроль в ОС Диспак демонстрируется на распределении бумаги: одному из отделов выделяется 100 м бумаги и формируется задача, заказывающая в паспорте заведомо больше распределенного количества м АЦПУ. Данная задача не будет сформирована. При заказе в паспорте метров АЦПУ меньше распределенного задача будет сформирована.\par
\par\textbf{5.2.10.} Для проверки ассемблера ВК «Эльбрус 1-КБ» с помощью программы «Димип» запускается задача «ТЕСТ3А» из библиотеки «КЗ» области «РЕЖИМ3».\par
\par\textbf{5.2.11.} Для проверки возможности решения задач в режиме 3 (расширенный формат) предназначена задача «Асскз», написанная на ассемблере ВК «Эльбрус 1-КБ». Задача запускается с помощью программы «Димип» из библиотеки «КЗ» области «РЕЖИМ3».\par
Процесс запуска задачи «Асскз» разбивается на три этапа:

\par\hangindent=1.5em - Формируется задача трансляции исходного текста. В результате трансляции на АЦПУ выдается листинг программы, а объектный модуль записывается в файл «Кзамод» указанной библиотеки;\par
\par\hangindent=1.5em - Формируется задача компоновки полученного объектного модуля. В результате компоновки на АЦПУ выдается листинг программы, а готовая программа записывается в файл «Кзагп»;\par
\par\hangindent=1.5em - Готовая программа запускается в решение. Результаты решения выдаются на АЦПУ.\par
\par\textbf{5.2.12.} Для проверки возможности вычисления элементарных функций в режиме 3 (расширенный формат) предназначена задача «Экстракод», написанная на ассемблере ВК «Эльбрус 1-КБ». Задача запускается с помощью программы «Димип» из библиотеки «КЗ» области «РЕЖИМ3».\par
Процесс запуска задачи «Экстракод» разбивается на три этапа:

\par\hangindent=1.5em - Формируется задача трансляции исходного текста. В результате трансляции на АЦПУ выдается листинг программы, а объектный модуль записывается в файл «Кзэмод» указанной библиотеки;\par
\par\hangindent=1.5em - Формируется задача компоновки полученного объектного модуля. В результате компоновки на АЦПУ выдается листинг программы, а готовая программа записывается в файл «Кзэгп»;\par
\par\hangindent=1.5em - Готовая программа запускается в решение. Результаты решения выдаются на АЦПУ.\par
\par\textbf{5.2.13.} Для проверки возможности решения задач на фортране в режиме 3 (транслятор FTN-В) предназначена задача «Кзфорт». Задача запускается с помощью программы «Димип» из библиотеки «КЗ» области «РЕЖИМ3».\par
Процесс запуска задачи «Кзфорт» разбивается на три этапа:

\par\hangindent=1.5em - Формируется задача трансляции исходного текста. В результате трансляции на АЦПУ выдается листинг программы, а объектный модуль записывается в файл «Кзфмод» указанной библиотеки;\par
\par\hangindent=1.5em - Формируется задача компоновки полученного объектного модуля. В результате компоновки на АЦПУ выдается листинг программы, а готовая программа записывается в файл «Кзфгп»;\par
\par\hangindent=1.5em - Готовая программа запускается в решение. Результаты решения выдаются на АЦПУ.\par
\par\textbf{5.2.14.} Проверяется возможность отладки программ, написанных на фортране, в адресном и символьном режимах. Задача подключается к отладчику «Слуга».\par
\par\textbf{5.2.15.} Статистическая информация о работе аппаратуры и прохождении задач выдается с помощью программ статистики «Сбор» и «Итог». Программы обрабативают информацию из области «Стат» и выдают результат на АЦПУ. Область «Стат» должна быть предварительно сформирована.\par
\par\textbf{5.2.16.} Проверка возможности формирования и накопления задач на буфере ввода и управления задачами осуществляется запуском пакета контрольных задач. Для запуска пакета контрольных задач после вызова терминальной задачи «Димип» требуется запустить в работу программу «Пакет N», пакет сформируется N раз. В случае отсутствия указания на число повторов пакет запустится только 1 раз.\par
Структура диалога для запуска пакета контрольных задач приведена в п.~6.16. каждая директива пользователя должна заканчиваться нажатием клавиши «Исполнение» на пульте терминала. Порядок формирования задач в пакете приведен в ТАБЛ.1.

\subsection*{5.3.~Проверка надежности}\addcontentsline{toc}{subsection}{5.3.~Проверка надежности}
\par\textbf{5.3.1.} Проверка надежности работы ОСПО осуществляется путем непрерывного двадцатичетырехчасового прогона пакета контрольных задач.\par
\par\textbf{5.3.2.} Для запуска пакета контрольных задач используется программа «Димип». Из библиотеки «Прогон» формируется задача «Прогон», которая многократно формирует контрольные задачи в количестве, обеспечивающем необходимую загрузку процессора. Описание контрольных задач приведено в приложении. Структура диалога для просмотра пакета контрольных задач и запуска их в решение приведена в п.~6.17.\par
\par\textbf{5.3.3.} Сформированный пакет задач производит полную проверку\par
Математического обеспечения, а непрерывность прогона и мультипрограммный режим работы производит эту проверку в различных сочетаниях.

\par\textbf{5.3.4.} Проверка идентичности (правильности) решения производится путем сличения результатов контрольного просчета с результатами, приведенными в приложении для каждой конкретной задачи.\par
\par\textbf{5.3.5.} По усмотрению комиссии в целях экономии бумаги на АЦПУ могут выдаваться только результаты счета контрольных задач (без информации о трансляции и загрузке).\par
\par\textbf{5.3.6.} На весь период испытаний распространяются следующие положения:\par
\par\hangindent=1.5em 1) При оценке результатов испытаний следует использовать следующие понятия:\par
\par\hangindent=1.5em - Отказ ОСПО --- Событие, состоящее в частичной или полной утрате работоспособности ОСПО и приводящее к стабильному невыполнению или неправильному выполнению задач под управлением ОС;\par
Для восстановления работоспособности ОСПО при отказе требуется перепрограммирование отдельных частей ОСПО;

\par\hangindent=1.5em - Сбой ОСПО --- Событие, состоящее во временной утрате работоспособности ОСПО и характеризуемое возникновением ошибки при выполнении задач под управлением ОС.\par
Для восстановления работоспособности ОСПО при сбое требуется проведение повторных действий по решению задачи или повторных вызовов операционной системы.

\par\hangindent=1.5em 2) Сбои ЭВМ, устраняемые программно-аппаратными средствами автоматически, при испытаниях не учитывают.\par
\par\hangindent=1.5em 3) Отказы и сбои ЭВМ, выявленные во время работ по техническому обслуживанию, не учитывают при определении количественных значений показателей надежности. Отказы зависимые (вторичные) не учитывают.\par
\par\hangindent=1.5em 4) В случае обнаружения ошибки в ОСПО испытание по данному пункту программы прекращается. Комиссия совместно с разработчиками устанавливает и анализирует причины, вызвавшие данную ошибку.\par
После устранения обнаруженной ошибки испытания по данному пункту повторяются.

Если, по мнению комиссии, обнаруженная ошибка может влиять на результаты ранее проведенных испытаний, испытания по данному пункту программы могут быть повторены.

\par\hangindent=1.5em 5) Сбои и отказы внешних устройств (МД, МЛ, АЦД, АЦПУ), не считаются сбоями ОСПО. Испытания продолжаются с использованием резервных внешних устройств.\par
В случае отказа устройства (пакета) с МД 2053 пакет задач должен быть введен повторно на резервном устройстве (пакете).

\par\textbf{5.3.7.} Для проверки ОСПО по отдельным пунктам программных испытаний члены комиссии могут представлять дополнительно свои пакеты программ, написанных на языках Фортран, Алгол-БЭСМ, ассемблер ВК «Эльбрус 1-КБ». Указанные программы должны удовлетворять следующим требованиям:\par
\par\hangindent=1.5em 1) Должны быть отлажены;\par
\par\hangindent=1.5em 2) Снабжены комплектом документации, содержащей:\par
\par\hangindent=1.5em - Перечень характеристик параметров ОСПО, проверяемых\par
С помощью данной программы;

\par\hangindent=1.5em - Аннотацию алгоритмов;\par
\par\hangindent=1.5em - Описание режимов функционирования и правил работы с ней;\par
\par\hangindent=1.5em - Текст программы на носителях,\par
Совместимых с носителями ВК «Эльбрус 1-КБ»;

\par\hangindent=1.5em - Листинг программы;\par
\par\hangindent=1.5em - Эталонные результаты решения задачи;\par
\par\hangindent=1.5em - Количественные характеристики решения задачи на других\par
Вычислительных средствах.

\par\textbf{5.3.8.} После окончания испытаний по каждому пункту «Методики» оформляется протокол. Допускается по усмотрению комиссии оформлять протокол по нескольким пунктам «Методики».\par
Результаты испытаний ОСПО в целом оформляются актом.

Настоящая «Методика» может корректироваться в процессе испытаний по требованию членов комиссии установленным порядком.

\section*{6.~Методы испытаний}\addcontentsline{toc}{section}{6.~Методы испытаний}
\par\textbf{6.1.} Предъявленная документация изучается на соответствие ведомости эд ИЫ.00110-01 20 01 и формуляру ИЫ.00110-01 30 01.\par
\begin{center}(п.~4.1, 4.2)\end{center}
\par\textbf{6.2.} Для демонстрации разметки магнитного диска необходимо поставить на накопитель гибких магнитных дисков (НГМД) №\,1 пультовой машины гибкий магнитный диск (ГМД) с программами разметки МД и вызова ОС. Вызвать программу разметки магнитного диска (см. «ОС Диспак. Руководство системного программиста» ИЫ.00110-01 32 01).\par
Размечается пакет №\,2053 (шестнадцатиричный код 805). После окончания разметки программа выходит на «Стоп 0» ( на РК код 0 3F 07 000).

\begin{center}(п.~3.2.11)\end{center}
\par\textbf{6.3.} Осуществляется перепись ОС с МЛ1 на вновь размеченный МД 2053 и производится вызов операционной системы Диспак. Согласно документу «ОС Диспак. Руководство системного программиста» ИЫ.00110-01 32 01.\par
По умолчанию задан МД 2053 (шестнадцатиричный код 805), канал 0, устройство 0.

После сообщения на операторский терминал №\,1 о ресурсах ОС с операторского терминала выполняются директивы начала работы системы: «Сме N», «Дат XXZZ», «Вре ЧЧММ».

Для ознакомления с составом ОСПО выдается сводка по архивным томам 2048 и 2053.

\begin{center}Структура диалога\end{center}
\begin{center}\small
\begin{BVerbatim}
              -----------------
Директивы пользователя       комментарии
----------------------       -----------
Сводка 419900002048      сводка по тому 2048
Сводка 419900002053      сводка по тому 2053
   На АЦПУ распечатывается структура архивных томов.
\end{BVerbatim}
\end{center}

\begin{center}(п.~3.1)\end{center}
\par\textbf{6.4.} По ключу арфа с шифром 419966 вызывается диалоговая программа «Арфа», с помощью которой создаются четыре области с именами КЗ31, КЗ32, КЗ33, КЗ34 --- Длиной L=200 зон каждая, чтение и запись разрешены только задачам с шифром 4199..\par
\begin{center}Структура диалога\end{center}
\begin{center}\small
\begin{BVerbatim}
          -----------------
Директива пользователя          ответ системы
----------------------          -------------
Арфа 419966                ***а R F а***версия
Нов КЗ31/4199..:L=200            выполнено
Нов КЗ32/4199..:L=200            выполнено
Нов КЗ33/4199..:L=200            выполнено
Нов КЗ34/4199..:L=200            выполнено
  КНЦ                           конец сеанса
\end{BVerbatim}
\end{center}

По усмотрению комиссии могут быть проверены другие директивы работы с архивной системой (см. «Операционная система Диспак. Руководство программиста.» ИЫ.00110-01 33 01, раздел 4).

\begin{center}(п.~3.2.8, 3.2.9)\end{center}
\par\textbf{6.5.} Проверка защиты информации от несанкционированного доступа в архиве осуществляется с помощью программы «Серп» (см. «Операционная система Диспак. Руководство программиста» ИЫ.00110-01 33 01, раздел 6). Программа «Серп» вызывается с шифром 1317, директивой «Проз» читается область «КЗ32», чтение из которой разрешено только задачам с шифром 4199... Предварительно с помощью программы «Арфа» демонстрируются права на область «КЗ32». При попытке чтения программа «Серп» снимается с диагностикой на АЦПУ «Чужая информация». Проверка по методике считается выполненной, если не удается прочитать информацию из чужой области.\par
\begin{center}Структура диалога\end{center}
\begin{center}\small
\begin{BVerbatim}
          -----------------
    Ответ системы         директива пользователя
    -------------         ----------------------
                                Арфа 419999
***ARFA***/ВЕРСИЯ/
 =                              кто КЗ32
<Сведения о правах на область>
                                КНЦ
 Конец сеанса
                                СП 1317
ИТМ и ВТ. С е р п:<версия>
 Что-                           про КЗ32/4199..
 Зона-                          10
 Конец сеанса
\end{BVerbatim}
\end{center}

\begin{center}(п.~3.2.9)\end{center}
\par\textbf{6.6.} Для проверки возможности вывода информации на внешние устройства в различных режимах одно из АЦПУ, например ацп 0, с помощью директивы «Вывод» программы «Серп» закрепляется за шифром 1317..... (см. «Операционная система Диспак. Руководство системного программиста.» ИЫ.00110-01 32 01).\par
До распределения АЦПУ выдачи поступали на свободное устройство. Вызывается тест АЦПУ с шифром 131755. Результат решения указанной задачи должен быть распечатан на распределенном АЦПУ (ацп 0). Выдачи задач с другими шифрами на данное АЦПУ не должны попадать. Распределение выдач на АЦПУ можно просмотреть директивой номв (см.»операционная система Диспак. Руководство оператора» ИЫ.00110-01 34 01).

\begin{center}Структура диалога\end{center}
\begin{center}\small
\begin{BVerbatim}
          -----------------
    Ответ системы         директива пользователя
    -------------         ----------------------
                              СП 419977
ИТМ и ВТ. С е р п:<версия>
Что-                          вывод
Пароль-                       02610
Шифр-                         1317.....
Класс-                        7
Устр-                         1А0
Есть
Шифр-                        <пустая посылка>
Маска.устр-                   1А0
Есть
Маска.устр-                   КНЦ
 Конец сеанса
                              Тестац 131755
Нвве
 Конец сеанса
\end{BVerbatim}
\end{center}

\begin{center}(п.~3.2.5, 3.2.7)\end{center}
\par\textbf{6.7.} Выключить оба АЦПУ. Директивой «Вывод» программы «Серп» снять распределение АЦПУ. Диалоговую работу вести с параллельной выдачей на АЦПУ. По окончании работы с программой «Серп» формируется выдача на АЦПУ. Так как АЦПУ выключено, на операторском терминале появится сообщение «Сбой АЦПУ N». Директива «Номв» подтвердит реакцию ОС на сбой устройства. После включения АЦПУ дается директива «Пов» и выдача продолжается. Проверяется реакция системы на сбои внешних устройств.\par
\begin{center}Структура диалога\end{center}
\begin{center}\small
\begin{BVerbatim}
          -----------------
    Ответ системы         директива пользователя
    -------------         ----------------------
                              СП 419977
ИТМ и ВТ. С е р п:<версия>
Что-                          а.
Уст а.
-                             Вывод
Пароль-                       02610
Шифр-                         убр 1317.....
Есть
Шифр-                        <пустая посылка>
Маска.устр-                   убр
Есть
Маска устр-                   КНЦ
Конец сеанса
\end{BVerbatim}
\end{center}

\begin{center}(п.~3.2.6)\end{center}
\par\textbf{6.8.} Для проверки бюджетно-статистического контроля в ОС с помощью программы «Бюджет», вызываемой по ключу «Бюджет» с шифром 4199, отделу 82 выделяется 100 м бумаги. С помощью программы «Димип» задаче МО5 в паспорте меняется шифр на «Шифр 820000000015ЗСХ» и заказывается 128 м бумаги, выполняются действия по формированию задачи МО5. В результате указанных действий задача не будет сформирована, на терминале появится сообщение об ошибках ввода из-за отсутствия метров АЦПУ. Директивой оператора «БДЖ 82» просматривается бюджет отдела 82. Затем в паспорте задачи заказывается 64 м АЦПУ и производятся действия по формированию задачи; задача формируется.\par
\begin{center}Структура диалога\end{center}
\begin{center}\small
\begin{BVerbatim}
          -----------------
    Ответ системы         директива пользователя
    -------------         ----------------------
                               Бюджет 4199
Кто-                           дис
Пароль-                        вахтер
Что-                           нов
Что-                           ВВД
Отдел                          82
-                              Б 100
Отдел                         <пустая посылка>
Записать-                     <пустая посылка>
Что-                           нбю
Что-                           пок 82
<Информация об отделе>
Что-                            КНЦ
 Конец сеанса
                               Димип 131700
=                              КТ прогон 0 прогон
=                              ред МО5
*                              л
<Текст задачи>                <редактирование>
*                              ф
<Информация об ошибках ввода - Нет метров АЦПУ>
=                              ред МО5
*                              л
<Текст задачи>                <редактирование>
*                              ф
*                              вый
 Конец сеанса
\end{BVerbatim}
\end{center}

\begin{center}(п.~3.2.4)\end{center}
\par\textbf{6.9.} Для проверки ассемблера ВК «Эльбрус 1-КБ» программой «Димип» формируется задача «ТЕСТ3А». Программа запуска задачи находится в файле «ПУСК3А». Результаты решения приведены в приложении.\par
\begin{center}Структура диалога\end{center}
\begin{center}\small
\begin{BVerbatim}
         -----------------
   Ответ системы         директива пользователя
   -------------         ----------------------
                          Димип 419900
=                         КТ РЕЖИМ3 0 КЗ
=                         ред ПУСК3А
*                         ф
                          Вый
Конец сеанса
\end{BVerbatim}
\end{center}

\begin{center}(п.~3.2.12)\end{center}
\par\textbf{6.10.} Для проверки работы ОС в режиме 3 (режим расширенного формата) с помощью программы «Димип» формируется задача «Асскз», написанная на ассемблере ВК «Эльбрус 1-КБ». Каждое следующее формирование задается только после окончания предыдущего этапа. Результаты решения приведены в приложении.\par
\begin{center}Структура диалога\end{center}
\begin{center}\small
\begin{BVerbatim}
          -----------------
    Ответ системы         директива пользователя
    -------------         ----------------------
                           Димип 419900
=                          КТ РЕЖИМ3 0 КЗ
=                          ред кзасс
*                          ф
*                          ред компа
*                          ф
*                          ред пуска
*                          ф
*                          вый
 Конец сеанса
\end{BVerbatim}
\end{center}

\begin{center}(п.~3.2.12)\end{center}
\par\textbf{6.11.} Для проверки вычисления элементарных функций в режиме 3 (режим расширенного формата) с помощью программы «Димип» формируется задача «Экстракод», написанная на ассемблере ВК «Эльбрус 1-КБ». Каждое следующее формирование задается только после окончания предыдущего этапа. Результаты решения приведены в приложении.\par
\begin{center}Структура диалога\end{center}
\begin{center}\small
\begin{BVerbatim}
          -----------------
    Ответ системы         директива пользователя
    -------------         ----------------------
                           Димип 419900
=                          КТ РЕЖИМ3 0 КЗ
=                          ред транэ
*                          ф
*                          ред компэ
*                          ф
*                          ред пускэ
*                          ф
*                          вый
 Конец сеанса
\end{BVerbatim}
\end{center}

\begin{center}(п.~3.2.12)\end{center}
\par\textbf{6.12.} Для проверки возможности решения задач на фортране в режиме 3 (режим расширенного формата) с помощью программы «Димип» формируется задача «Кзфорт», написанная на фортране (транслятор FTN-В). Каждое следующее формирование задается только после окончания предыдущего этапа. Результаты решения приведены в приложении.\par
\begin{center}Структура диалога\end{center}
\begin{center}\small
\begin{BVerbatim}
          -----------------
    Ответ системы         директива пользователя
    -------------         ----------------------
                          Димип 419900
=                         КТ РЕЖИМ3 0 КЗ
=                         ред кзфорт
*                         ф
*                         ред компф
*                         ф
*                         ред пускф
*                         ф
*                         вый
 Конец сеанса
\end{BVerbatim}
\end{center}

\begin{center}(п.~3.2.13)\end{center}
\par\textbf{6.13.} Возможность интерактивной отладки программ проверяется запуском задачи МО5, написанной на языке Фортран, с разделом теле и с управляющей картой *DEBUG в функциональной части пакета для последующей постановки ее на диалоговый отладчик «Слуга», в соответствии с нижеприведенной структурой диалога.\par
\begin{center}Структура диалога\end{center}
\begin{center}\small
\begin{BVerbatim}
              -----------------
  Директива                ответ     комментарии
  ---------               -------    -----------
Димип 419999                =
КТ прогон 0 прогон          =
Ред МО5                     *       вызов задачи МО5
З 1                        00001
\end{BVerbatim}
\end{center}

Шифр 419905ЗС2‾ТЕЛ‾ОЗУ 32‾ 00002 введение раздела теле

<Исполнение> *

В 4 00005

*DEBUG 00006

<Исполнение> * введение у/к *DEBUG

Ф * формирование задачи М05

Вый конец сеанса окончание работы

\begin{center}С программой «Димип»\end{center}
Слуга слуга(версия) вызов программы

\begin{center}Шифр «Слуга»\end{center}
419905 есть -

Пва уст. Пва установка режима

V уст. V печати протокола

\begin{center}И выдачи на экран.\end{center}
РЖ а адресный режим.

К 1000 --- Установка контрольной

\begin{center}Точки по адресу 1000\end{center}
\begin{center}В мониторной системе\end{center}
И 01000 кра останов в задаче по

\begin{center}Адресу 1000.\end{center}
У 1003 --- Установка узла

И 01003 узел останов в задаче по

\begin{center}Узлу\end{center}
РЖ ф --- Установка режима

\begin{center}Отладки на фортране\end{center}
\begin{center}(Символьный режим)\end{center}
Жду PROGRAM --- Останов в задаче по

И раздел PROGRAM передаче управления

\begin{center}(Ф):тест на головной раздел\end{center}
ЗП М3:DOLOOP --- Установка КТ по запи-

\begin{center}Си в ячейку с иденти-\end{center}
\begin{center}Фикатором М3 в под-\end{center}
\begin{center}Программе DOLOOP\end{center}
И (ф)оператор 7: останов по записи в

\begin{center}DOLOOP ячейку с идентифика-\end{center}
\begin{center}01047П по зап. Тором М3 в подпрог-\end{center}
\begin{center}Рамме DOLOOP\end{center}
М2:DOLOOP= 6400000000000002 выдача содержимого

\begin{center}Ячейки с идентифика-\end{center}
\begin{center}Тором М2 в подпрог-\end{center}
\begin{center}Рамме DOLOOP в виде\end{center}
\begin{center}Логической константы\end{center}
М2:DOLOOP:=

6400000000000003 --- Изменение значения

\begin{center}Ячейки с идентифика-\end{center}
\begin{center}Тором М2 в подпрог-\end{center}
\begin{center}Рамме DOLOOP\end{center}
М2:DOLOOP=Д 3000000000+01 выдача значения

\begin{center}Ячейки с идентифика-\end{center}
\begin{center}Тором М3 в подпрог-\end{center}
\begin{center}Рамме DOLOOP в виде\end{center}
\begin{center}Десятичного нормали-\end{center}
\begin{center}Зованного числа\end{center}
М2:DOLOOP:=

6400000000000002 --- Изменение значения

\begin{center}Ячейки с идентифика-\end{center}
\begin{center}Тором М2 в подпрог-\end{center}
\begin{center}Рамме DOLOOP\end{center}
К 66:DOLOOP --- Установка контрольной

\begin{center}Точки на операторе 66\end{center}
\begin{center}В подпрограмме DOLOOP\end{center}
\begin{center}По кра\end{center}
И выдача значений сумматора, регистра младших

\begin{center}Разрядов, регистра команд, индексных регист-\end{center}
\begin{center}Ров, регистра 16 задачи и исполняемой команды\end{center}
К 70:DOLOOP установка контрольной

\begin{center}Точки на операторе 70\end{center}
\begin{center}В подпрограмме DOLOOP\end{center}
\begin{center}По кра\end{center}
И (ф) оператор 70:DOLOOP

\begin{center}01547 по кра останов по кра на опе-\end{center}
\begin{center}Раторе 70 в подпрог-\end{center}
\begin{center}Рамме DOLOOP\end{center}
И (Ф)13:STOP*

\begin{center}01027 конец задачи конец задачи\end{center}
КНЦ конец сеанса окончание работы с

\begin{center}Отладчиком\end{center}
\begin{center}(п.~3.2.10)\end{center}
\par\textbf{6.14.} Для проверки возможности планирования мультипрограммной смеси выдается директива «Дай спе» и вызывается диалоговая программа «Димип». Однократно формируется пакет контрольных задач. Проверяется выполнение директив управления и контроля за ходом решения задач согласно «Руководству оператора» ИЫ.00110-01 34 01.\par
Проверяются директивы информационные --- «ВЦПП», «Ном», «СКВ», «СКВ <номер на буфере ввода>», «Пак»; директивы управления --- «ВКН <номер канала>», «Выб <номер на буфере ввода>», «Нет спе» и ДР. По усмотрению комиссии.

\begin{center}Структура диалога\end{center}
\begin{center}\small
\begin{BVerbatim}
         -----------------
   Ответ системы         директива пользователя
   -------------         ----------------------
                         Димип 419944
=                        КТ прогон 0 прогон
=                        пакет
БВ <N СКВ> задача МО2
БВ <N СКВ> задача КЗ1
\end{BVerbatim}
\end{center}

<Формирование задач на буфер ввода ОС, N-НОМЕР СКВ>

Прошли круг 1

\begin{center}Вый\end{center}
Конец сеанса

\begin{center}(п.~3.2.3)\end{center}
\par\textbf{6.15.} Для выдачи статистической информации о работе аппаратуры и пользователей с любого терминала вызывается программа «Сбор» по ключу «Сбстат 415501». После окончания программы «Сбор» на АЦПУ выдается распечатка о работе аппаратуры. Информацию о задачах пользователя выдает программа «Итог». Она вызывается с помощью программы «Димип» из области «Пускст».\par
\begin{center}Структура диалога\end{center}
\begin{center}\small
\begin{BVerbatim}
         -----------------
   Ответ системы         директива пользователя
   -------------         ----------------------
                         Димип 415544
=                        КТ пускст 0 стат
=                        ред пускст
*                        л
<Текст файла>
                         Ф 3
                         Вый
Конец сеанса
\end{BVerbatim}
\end{center}

\begin{center}(п.~3.2.4)\end{center}
\par\textbf{6.16.} Для проверки возможности формирования и накопления задач на буфере ввода выдается директива «Дай спе» и формируется пакет контрольных задач 3 раза. Накопление задач в очереди проверяется директивами «СКВ», «Ном», «Ном N» и ДР. (см. «Операционная система Диспак. Руководство оператора» ИЫ.00110-01 34 01).\par
\begin{center}Структура диалога\end{center}
\begin{center}\small
\begin{BVerbatim}
          -----------------
    Ответ системы         директива пользователя
    -------------         ----------------------
                           Димип 419999
 =                         КТ прогон 0 прогон
 =                         пакет 3
 БВ <N СКВ> задача МО2
 БВ <N СКВ> задача КЗ1
 <Формирование задач на буфер ввода ОС>
 Прошли круг 1
 БВ <N СКВ> задача МО2
 БВ <N СКВ> задача КЗ1
<Формирование задач на буфер ввода ОС>
 Прошли круг N  <N - Параметр задания пакет N>
                           Вый
 Конец сеанса
\end{BVerbatim}
\end{center}

\begin{center}(п.~3.2.2)\end{center}
\par\textbf{6.17.} Проверка и оценка возможности круглосуточной непрерывной мультипрограммной работы ОСПО производится путем непрерывного круглосуточного прогона пакета контрольных задач. Набор задач рассчитан на непрерывный суточный прогон. Для формирования пакета задач предназначена программа «Прогон», которая вызывается из библиотеки «Прогон» с помощью программы «Димип».\par
\begin{center}Структура диалога\end{center}
\begin{center}\small
\begin{BVerbatim}
         -----------------
   Ответ системы         директива пользователя
   -------------         ----------------------
                          Димип 419999
=                         КТ прогон 0 прогон
=                         ред прогон
*                         л
<Текст программы>
                          Прогон
<Формирование задач на буфер ввода ОС>
БВ <N СКВ> задача МО1
БВ <N СКВ> задача КЗ5
Все                       вый
Конец сеанса
\end{BVerbatim}
\end{center}

Названия контрольных задач приведены в диалоге в качестве примера. Результаты решения контрольных задач должны совпадать с эталонными результатами, приведенными в приложении.

\begin{center}(п.~3.2.1)\end{center}
\par\textbf{6.18.} ОСПО считается выдержавшим испытания, если:\par
\par\hangindent=1.5em - Выполнены все требования раздела 3;\par
\par\hangindent=1.5em - В течение прогона не было ни одного отказа ОСПО, а число сбоев не превысило четырех;\par
\par\hangindent=1.5em - Эксплуатационная документация обеспечивает полное использование возможностей, предоставляемых ОСПО пользователю (программисту, оператору).\par
По результатам испытаний составляется акт приемо-сдаточных испытаний. Примечания: 1. В скобках приведены пункты технических требований, проверяемые данной методикой.

2. В ТАБЛ.2 приведено соответствие пунктов проверяемых технических требований и пунктов методики проверки.

\begin{flushright}\small Таблица 1.\end{flushright}
\begin{center}Порядок ввода задач\end{center}
\begin{center}\small\begin{tabularx}{\linewidth}{|>{\raggedright\arraybackslash\hsize=0.516\hsize}X|>{\raggedright\arraybackslash\hsize=0.710\hsize}X|>{\raggedright\arraybackslash\hsize=1.548\hsize}X|>{\raggedright\arraybackslash\hsize=1.226\hsize}X|}\hline
\small\bfseries\hspace{0pt}N П/П & \small\bfseries\hspace{0pt}N задачи & \small\bfseries\hspace{0pt}Язык программирования & \small\bfseries\hspace{0pt}Шифр задач \\ \hline
1 & МО2 & автокод ИТМ & 419900000012 \\ \hline
2 & КЗ 1 & Машинные коды & 419900000001 \\ \hline
3 & Мо 5 & Фортран & 419900000015 \\ \hline
4 & КЗ 5 & Машинные коды & 419900000005 \\ \hline
5 & Мо 1 & Фортран & 419900000011 \\ \hline
6 & КЗ 3 & Машинные коды & 419900000003 \\ \hline
7 & Мо 4 & Фортран & 419900000014 \\ \hline
8 & КЗ 2 & Машинные коды & 419900000002 \\ \hline
9 & Мо 3 & Фортран & 419900000013 \\ \hline
10 & КЗ 5 & Машинные коды & 419900000005 \\ \hline
11 & Мо 9 & Алгол-БЭСМ & 419900000019 \\ \hline
12 & КЗ 2 & Машинные коды & 419900000002 \\ \hline
13 & КЗ 3 & Машинные коды & 419900000003 \\ \hline
14 & Мо 1 & Фортран & 419900000011 \\ \hline
15 & Мо 3 & Фортран & 419900000013 \\ \hline
16 & КЗ 5 & Машинные коды & 419900000005 \\ \hline
17 & Мо 4 & Фортран & 419900000014 \\ \hline
18 & Мо 1 & Фортран & 419900000011 \\ \hline
\end{tabularx}\end{center}

\begin{flushright}\small Таблица 2.\end{flushright}
\begin{center}\small\begin{tabularx}{\linewidth}{|>{\raggedright\arraybackslash\hsize=1.040\hsize}X|>{\raggedright\arraybackslash\hsize=0.960\hsize}X|}\hline
\small\bfseries\hspace{0pt}Пункт требований & \small\bfseries\hspace{0pt}Пункт методики \\ \hline
3.1 & 6.3. \\ \hline
3.2.1 & 6.17 \\ \hline
3.2.2. & 6.16 \\ \hline
3.2.3 & 6.14 \\ \hline
3.2.4 & 6.8,6.15 \\ \hline
3.2.5 & 6.6 \\ \hline
3.2.6. & 6.7 \\ \hline
3.2.7 & 6.6 \\ \hline
3.2.8 & 6.4 \\ \hline
3.2.9 & 6.4, 6.5 \\ \hline
3.2.10 & 6.13 \\ \hline
3.2.11 & 6.2 \\ \hline
3.2.12 & 6.9-6.11 \\ \hline
3.2.13 & 6.12 \\ \hline
4.1, 4.2 & 6.1 \\ \hline
\end{tabularx}\end{center}

\section*{Приложение}\addcontentsline{toc}{section}{Приложение}
\begin{center}Описание контрольных задач\end{center}
\section*{1.~Контрольная задача КЗ1}\addcontentsline{toc}{section}{1.~Контрольная задача КЗ1}
\begin{center}\small
\begin{BVerbatim}
       -------------------------
 Решение уравнений внешней баллистики
Методом рунге-кутта с переменным шагом
\end{BVerbatim}
\end{center}

\par\textbf{1.1.} Программа предназначена для решения 18-ТИ уравнений внешней баллистики с автоматическим выбором шага (начальный шаг задается). На каждом шаге интегрирования производится двойной просчет. На следующий шаг интегрирования идем только в том случае, если данные, полученные при двух последовательных просчетах, совпадают.\par
\subsection*{1.2.~Шифр задачи 419900000001}\addcontentsline{toc}{subsection}{1.2.~Шифр задачи 419900000001}
Программа занимает 2 листа оперативной памяти;

Решается в первом режиме.

Время решения --- 3.50 мин.

\par\textbf{1.3.} Результаты выдачи: 1-ОЕ число в 1-ОЙ строке: +232830643653Е-09, 1-ОЕ число в 4-ОЙ строке: -245623186191Е+05, последнее число в 4-ОЙ строке: -107218410769Е-01\par
\section*{2.~Контрольная задача КЗ2}\addcontentsline{toc}{section}{2.~Контрольная задача КЗ2}
\begin{center}\small
\begin{BVerbatim}
------------------------
\end{BVerbatim}
\end{center}

Плоская задача неустановившихся движений тяжелой жидкости

\par\textbf{2.1.} Программа предназначена для проверки функционирования операционной системы.\par
Задача идет с двойным просчетом. При несовпадении повторного счета на печать выводятся две контрольные суммы и счет продолжается. В конце решения задача выдает на АЦПУ распечатку, состоящую из 200 выдач.

\subsection*{2.2.~Шифр задачи 419900000002}\addcontentsline{toc}{subsection}{2.2.~Шифр задачи 419900000002}
Программа занимает 8 листов оперативной памяти;

Решается в первом режиме;

Задаче требуется 32 тракта;

Время решения --- 5.15 мин.

\par\textbf{2.3.} Результаты выдачи:\par
Выдача 1: 1-ОЕ число последней строки +2795188471Е-04

Выдача 100: 1-ОЕ число последней строки +1252317649Е+00

Выдача 200: 1-ОЕ число последней строки +2736036693Е+00.

\section*{3.~Контрольная задача КЗ3}\addcontentsline{toc}{section}{3.~Контрольная задача КЗ3}
\begin{center}\small
\begin{BVerbatim}
     ------------------------
Обращение матрицы большого порядка
\end{BVerbatim}
\end{center}

\par\textbf{3.1.} Программа использует модифицированную схему жордана. При обращении матрицы большого порядка, когда она не помещается в память, разбиваем ее на клетки. Размерность клетки 3232, что соответствует объему одного листа оперативной памяти или одной зоне МД. Программа предназначена для проверки функционирования системы при большом количестве обменов с магнитными дисками. В процессе решения информация записывается в 4 области.\par
\subsection*{3.2.~Шифр задачи 419900000003}\addcontentsline{toc}{subsection}{3.2.~Шифр задачи 419900000003}
Программа занимает 5 листов оперативной памяти;

Задаче требуется 4 области на магнитных дисках

КЗ31, КЗ32, КЗ33, КЗ34, доступные для записи;

Решается в первом режиме.

Время решения --- 0.38 мин.

\par\textbf{3.3.} Результаты выдачи:\par
\begin{center}1 строка: +6400001876Е+03\end{center}
Последняя строка: +9599998875Е+03

\section*{4.~Контрольная задача КЗ5}\addcontentsline{toc}{section}{4.~Контрольная задача КЗ5}
\begin{center}\small
\begin{BVerbatim}
       -------------------------
Расчет энергии  взаимодействия молекул
       Кристаллического вещества
\end{BVerbatim}
\end{center}

\par\textbf{4.1.} Программа предназначена для расчета энергии и плотности энергии взаимодействия молекул кристаллического вещества и производных до 2-ГО порядка от этих функций по параметрам состояния вещества.\par
Для контроля вычислений используется система двойного просчета. В результате решения на АЦПУ поступает 30 выдач из 26-ТИ строк по 8 чисел в каждой. На печать выдается значение функций и их производных. В случае несовпадения контрольных сумм 2-Х просчетов на печать выдаются также контрольные суммы. В случае 40-КРАТНОГО повторения одного просчета и несовпадения никаких двух сумм происходит окончание задачи по адресу 02055. Контрольные суммы выдаются в виде констант.

\subsection*{4.2.~Шифр задачи 419900000005}\addcontentsline{toc}{subsection}{4.2.~Шифр задачи 419900000005}
Программа занимает 4 листа оперативной памяти;

Задаче требуется 32 тракта;

Решается в первом режиме;

Время решения --- 21.10 мин.

\par\textbf{4.3.} Результаты выдачи:\par
1 выдача: 1-ОЕ число последней строки: +3462015909Е-01 15 выдача: 1-ОЕ число последней строки: +3429745207Е-01 30 выдача: 1-ОЕ число последней строки: +3461995815Е-01

\section*{5.~Контрольная задача МО1}\addcontentsline{toc}{section}{5.~Контрольная задача МО1}
\begin{center}\small
\begin{BVerbatim}
------------------------
\end{BVerbatim}
\end{center}

\par\textbf{5.1.} Проверяется транслятор с языка Фортран, версия компилятора «FOREX». В задаче определяется плотность частиц пучка. Идентификатор задачи --- «Контрольная задача». В результате трансляции и счета задача выдает результаты решения на АЦПУ.\par
\subsection*{5.2.~Шифр задачи 419900000011}\addcontentsline{toc}{subsection}{5.2.~Шифр задачи 419900000011}
Задача занимает 26 листов оперативной памяти;

Решается в первом режиме;

Время трансляции и решения --- 11.57 мин.

\par\textbf{5.3.} Результаты выдачи на АЦПУ:\par
\begin{center}TIME=20.0\end{center}
Х DENSITY U

0.00 0.1336 0.0000

2.00 1.3914 1.2610

4.00 2.4650 5.3531

18.00 2.9353 45.0905

98.00 2.9909 283.3802

100.00 2.9911 289.3621

\begin{center}TIME=40.0\end{center}
Х DENSITY U

0.00 0.0059 0.0000

2.00 0.1187 0.0795

100.00 2.9818 278.7566

\section*{6.~Контрольная задача МО2}\addcontentsline{toc}{section}{6.~Контрольная задача МО2}
\begin{center}\small
\begin{BVerbatim}
-------------------------
\end{BVerbatim}
\end{center}

\par\textbf{6.1.} Программа служит для проверки автокода ИТМ. Ищется решение системы обыкновенных линейных уравнений 2-ГО порядка модифицированным методом прогонки.\par
\subsection*{6.2.~Шифр задачи 419900000012}\addcontentsline{toc}{subsection}{6.2.~Шифр задачи 419900000012}
Задача решается в первом режиме;

Время трансляции и решения --- 00.02 мин

\par\textbf{6.3.} Результаты выдачи:\par
3-Я строка: А2М=М, А1=N, А2Р=Р, ДИ1=Р(321)

Предпоследняя строка (в строке 10 результатов):

+7757Е-01 +6983Е-01 ...

\section*{7.~Контрольная задача МО3}\addcontentsline{toc}{section}{7.~Контрольная задача МО3}
\begin{center}\small
\begin{BVerbatim}
-------------------------
\end{BVerbatim}
\end{center}

\par\textbf{7.1.} Программа представляет собой контрольную задачу МО1, версия компилятора Фортран-ГДР.\par
\subsection*{7.2.~Шифр задачи 419900000013}\addcontentsline{toc}{subsection}{7.2.~Шифр задачи 419900000013}
Задача занимает 23 листа оперативной памяти;

Решается в первом режиме;

Время трансляции и решения --- 13.06 мин.

\subsection*{7.3.~Результаты выдачи см. задачу МО1}\addcontentsline{toc}{subsection}{7.3.~Результаты выдачи см. задачу МО1}
\section*{8.~Контрольная задача МО4}\addcontentsline{toc}{section}{8.~Контрольная задача МО4}
\begin{center}\small
\begin{BVerbatim}
-------------------------
\end{BVerbatim}
\end{center}

\par\textbf{8.1.} Программа представляет собой контрольную задачу МО1, версия компилятора Фортран-Дубна.\par
\subsection*{8.2.~Шифр задачи 419900000014}\addcontentsline{toc}{subsection}{8.2.~Шифр задачи 419900000014}
Задача занимает 23 листа оперативной памяти;

Решается в первом режиме;

Время трансляции и решения --- 13.58 мин.

\subsection*{8.3.~Результаты выдачи см. задачу МО1}\addcontentsline{toc}{subsection}{8.3.~Результаты выдачи см. задачу МО1}
\section*{9.~Контрольная задача МО5}\addcontentsline{toc}{section}{9.~Контрольная задача МО5}
\begin{center}\small
\begin{BVerbatim}
-------------------------
\end{BVerbatim}
\end{center}

\par\textbf{9.1.} Задача представляет собой тест проверки трансляции циклов. Проверяется система Фортран-Дубна.\par
\subsection*{9.2.~Шифр задачи 419900000015}\addcontentsline{toc}{subsection}{9.2.~Шифр задачи 419900000015}
Задача занимает 23 листа оперативной памяти;

Решается в первом режиме;

Время трансляции и решения --- 00.08 мин.

\par\textbf{9.3.} Результаты выдачи: последняя строка выдачи: ALL 22 CNECKS ARE окау\par
\section*{10.~Контрольная задача МО9}\addcontentsline{toc}{section}{10.~Контрольная задача МО9}
\begin{center}\small
\begin{BVerbatim}
--------------------------
\end{BVerbatim}
\end{center}

\par\textbf{10.1.} Решается задача об автомате G в стационарной среде. Проверяется система Алгол-БЭСМ. в результате трансляции и решения печатаются 99 строк.\par
\subsection*{10.2.~Шифр задачи 419900000019}\addcontentsline{toc}{subsection}{10.2.~Шифр задачи 419900000019}
Задача занимает 19 листов оперативной памяти;

Решается в первом режиме;

Время трансляции и решения --- 00.04 мин.

\par\textbf{10.3.} Результаты выдачи: последняя строка выдачи: 320.000000\par
\section*{11.~Контрольная задача «ТЕСТ3А»}\addcontentsline{toc}{section}{11.~Контрольная задача «ТЕСТ3А»}
\begin{center}\small
\begin{BVerbatim}
-------------------------------
\end{BVerbatim}
\end{center}

\par\textbf{11.1.} Задача предназначена для проверки ассемблера ВК «Эльбрус 1-КБ» третьего режима.\par
\subsection*{11.2.~Шифр задачи: 419930000000}\addcontentsline{toc}{subsection}{11.2.~Шифр задачи: 419930000000}
Задача занимает 32 листа оперативной памяти,

Решается в первом режиме;

Время трансляции и решения --- 00.13 мин.

\par\textbf{11.3.} Результаты решения:\par
В результате компиляции выдается полный листинг ассемблера. На последней странице листинга содержится текст:

Число ошибок 78, предупреждений 35

\section*{12.~Контрольная задача «Асскз»}\addcontentsline{toc}{section}{12.~Контрольная задача «Асскз»}
\begin{center}\small
\begin{BVerbatim}
--------------------------
\end{BVerbatim}
\end{center}

\par\textbf{12.1.} Контрольная задача «Асскз» демонстрирует основные возможности режима 3 ВК «Эльбрус 1-КБ» --- Использование расширенного адресного пространства, компонентных операций и группового режима.\par
В процессе счета генерируется квадратная матрица ранга 1000 специального вида. Затем полученная матрица траспонируется и обе матрицы перемножаются. Вычисляется L-НОРМА итоговой матрицы.

Задача написана на ассемблере ВК «Эльбрус 1-КБ» и выполняется в три этапа: трансляция, компоновка, решение готовой программы.

\subsection*{12.2.~Шифр задачи трансляции: 419930000001}\addcontentsline{toc}{subsection}{12.2.~Шифр задачи трансляции: 419930000001}
\begin{center}Шифр задачи компоновки: 419930000002\end{center}
\begin{center}Шифр задачи решения ГП: 419930000003\end{center}
Задача использует 3000К слов оперативной памяти;

Решается в третьем режиме;

Время решения --- 12.09 мин.

\par\textbf{12.3.} По окончании счета на АЦПУ выдаются несколько строк заголовка и печатается число:\par
+899032184389Е+19

\section*{13.~Контрольная задача «Экстракод»}\addcontentsline{toc}{section}{13.~Контрольная задача «Экстракод»}
\begin{center}\small
\begin{BVerbatim}
------------------------------
\end{BVerbatim}
\end{center}

\subsection*{13.1.~Вычисление стандартных функций}\addcontentsline{toc}{subsection}{13.1.~Вычисление стандартных функций}
Тест предназначен для проверки возможности использования комплекса программ вычисления стандартных функций. Программа написана на языке ассемблер ВК «Эльбрус 1-КБ». В тесте вычисляются значения:

SQRT(Х) для Х=5,2,3,10, пи

LN(Х) для Х=5,2,3,10, пи

LG(Х) для Х=5,2,3,10, пи

Ехр(х) для Х=1,3,2,4,10

SIN(Х) для Х=1, ПИ/4, ПИ/2, ПИ,1

COS(Х) для Х=1, ПИ/4, ПИ/2, ПИ,1. Цикл повторяется 1000 раз.

\par\textbf{13.2.} Прохождение задачи на ВК «Эльбрус 1-КБ» включает 3 этапа:\par
Трансляция, шифр 41993000011

Компоновка, шифр 41993000012

Пуск ГП, шифр 41993000013

Задача решается в третьем режиме;

Время решения: 00.01 мин.

\subsection*{13.3.~Результаты выдачи}\addcontentsline{toc}{subsection}{13.3.~Результаты выдачи}
1 число 1 строки +223606797749Е+01 (SQRT(5))

2 число 2 строки +693147180559Е+00 (LN(2) )

2 число 3 строки +301029995663Е+00 (LG(2) )

3 число 4 строки +738905609893Е+01 (ЕХР(2) )

2 число 5 строки +707106781185Е+00 (SIN(ПИ/4))

4 число 6 строки -1000000000000+01 (COS(ПИ))

\section*{14.~Контрольная задача «Кзфорт»}\addcontentsline{toc}{section}{14.~Контрольная задача «Кзфорт»}
\begin{center}\small
\begin{BVerbatim}
---------------------------
\end{BVerbatim}
\end{center}

\par\textbf{14.1.} Контрольная задача «Кзфорт» написана на фортране ВК «Эльбрус 1-КБ» (транслятор FTN-В). Массив большой длины (1000000 слов) заполняется псевдослучайными числами; затем массив упорядочивается методом быстрой сортировки. Задача выполняется в три этапа: трансляция, компоновка, решение готовой программы.\par
\subsection*{14.2.~Шифр задачи трансляции: 419931000001}\addcontentsline{toc}{subsection}{14.2.~Шифр задачи трансляции: 419931000001}
\begin{center}Шифр задачи компоновки: 419931000002\end{center}
\begin{center}Шифр задачи решения ГП: 419931000003\end{center}
Задача использует 1000К слов оперативной памяти;

Решается в третьем режиме;

Время решения --- 1.27 мин.

\par\textbf{14.3.} По окончании счета на АЦПУ выдаются несколько строк заголовка и печатаются первые 10 элементов упорядоченного по возрастанию массива:\par
Первое число первой строки: -989999262399Е+02 последнее число второй строки: -989964919396Е+02

\section*{15.~Контрольная задача «КЗ10»}\addcontentsline{toc}{section}{15.~Контрольная задача «КЗ10»}
\begin{center}\small
\begin{BVerbatim}
------------------------------
\end{BVerbatim}
\end{center}

\par\textbf{15.1.} Задача «КЗ10» написана на языке Фортран в мониторной системе «Дубна». Выполняется интегрирование системы дифференциальных уравнений. Программа в оттранслированном виде хранится в личной библиотеке.о 1253\par
\subsection*{15.2.~Шифр задачи: 41990000010}\addcontentsline{toc}{subsection}{15.2.~Шифр задачи: 41990000010}
Используемые ресурсы: 23 листа оперативной памяти;

11 трактов.

Задача решается в первом режиме.

Время РЕШЕНИЯ:20.26 мин.

\subsection*{15.3.~Результаты выдачи}\addcontentsline{toc}{subsection}{15.3.~Результаты выдачи}
Первая строка: ТНК, РН, РК

Вторая строка: 0.00000028800000.000000

1 число 3 строки: 4828.2382390

1 число последней строки: -0.2796701

\section*{16.~Контрольная задача КЗ11}\addcontentsline{toc}{section}{16.~Контрольная задача КЗ11}
\begin{center}\small
\begin{BVerbatim}
-----------------------
   Пакетный «UNICOP»
\end{BVerbatim}
\end{center}

\par\textbf{16.1.} Задача предназначена для проверки работы операционной системы Диспак с магнитными лентами типа ЕС и состоит из двух частей:\par
В первой части задачи с помощью пакетной программы «UNICOP» директивой «Б6ЕС» из области «Прогон» переписывается 100 зон на магнитную ленту 6. Для надежности создается два одинаковых файла --- Основной и дубль. На печать выдается контрольная сумма одного файла;

Во второй части задачи с помощью той же программы по директиве «Ескс» подсчитывается контрольная сумма информации на ленте. Вычисленная контрольная сумма сравнивается с суммой, полученной при записи. В случае несовпадения на печать выдаются номера записей с несовпавшими суммами, а также считанная и эталонная контрольные суммы.

\subsection*{16.2.~Шифр задачи 419900000111}\addcontentsline{toc}{subsection}{16.2.~Шифр задачи 419900000111}
Задаче требуется магнитная лента в режиме ЕС, предварительно ленте должен быть присвоен номер директивой оператора «Есм 6 НУ», где НУ --- Номер магнитофона с МЛ.

Решается в первом режиме.

Время решения --- 0.02 мин.

\par\textbf{16.3.} Результаты решения:\par
Визуально сравниваются контрольные суммы, полученные в каждой части задачи. Контрольные суммы должны быть равны между собой.

\section*{Перечень ссылочных документов}\addcontentsline{toc}{section}{Перечень ссылочных документов}
\section*{1.~Общее системное программное обеспечение}\addcontentsline{toc}{section}{1.~Общее системное программное обеспечение}
Спецификация.

ИЫ.00110-01. 2. Общее системное прогаммное обеспечение.

Операционная система Диспак.

Руководство системного программиста.

ИЫ.00110-01 32 01. 3. Общее системное прогаммное обеспечение.

Операционная система Диспак. Руководство программиста.

ИЫ.00110-01 33 01. 4. Общее системное прогаммное обеспечение.

Операционная система Диспак. Руководство оператора.

ИЫ.00110-01 34 01. 5. Система тестовых программ.

Руководство оператора.

ИЫ.04301-01 34 01. 6. Пультовое программное обеспечение.

Руководство оператора.

ИЫ.00112-01 34 01. 7. Вычислительный комплекс «Эльбрус 1-КБ»

Техническое описание.

ИЫ3.055.045 ТО5.

\section*{Лист регистрации изменений}\addcontentsline{toc}{section}{Лист регистрации изменений}
.-----------------------------------------------------------------.

\begin{center}\small\begin{tabularx}{\linewidth}{|>{\raggedright\arraybackslash\hsize=0.176\hsize}X|>{\raggedright\arraybackslash\hsize=0.706\hsize}X|>{\raggedright\arraybackslash\hsize=1.059\hsize}X|>{\raggedright\arraybackslash\hsize=1.059\hsize}X|>{\raggedright\arraybackslash\hsize=1.059\hsize}X|>{\raggedright\arraybackslash\hsize=1.059\hsize}X|>{\raggedright\arraybackslash\hsize=1.235\hsize}X|>{\raggedright\arraybackslash\hsize=1.588\hsize}X|>{\raggedright\arraybackslash\hsize=1.941\hsize}X|>{\raggedright\arraybackslash\hsize=1.235\hsize}X|>{\raggedright\arraybackslash\hsize=0.706\hsize}X|>{\raggedright\arraybackslash\hsize=0.176\hsize}X|}\hline
 & \small\bfseries\hspace{0pt}И з м. & \small\bfseries\hspace{0pt}Номер измененных & \small\bfseries\hspace{0pt}Лист замененных & \small\bfseries\hspace{0pt}В (СТ новых & \small\bfseries\hspace{0pt}Аниц) аннулированных & \small\bfseries\hspace{0pt}Всего листов (СТР.) В документе & \small\bfseries\hspace{0pt}Номер документа & \small\bfseries\hspace{0pt}Входящий номер сопроводи-тельного документа & \small\bfseries\hspace{0pt}Подпись & \small\bfseries\hspace{0pt}Д а т а &  \\ \hline
\end{tabularx}\end{center}

\end{document}
